\section{Vulnerability Engine Design}
\label{sec:engine}

The new \path{engine/} package is intended to become the framework's central
vulnerability-finding layer. It consists of three primary components:
\texttt{VulnScanner}, \texttt{BatchScanner}, and \texttt{FindingRouter}.

\begin{figure}[t]
  \centering
  \resizebox{0.95\textwidth}{!}{\begin{tikzpicture}[
  font=\scriptsize,
  stage/.style={rounded corners=3pt, draw=raBlue, fill=raBlue!7, thick, minimum width=0.17\textwidth, text width=0.17\textwidth, minimum height=1.25cm, align=center},
  metric/.style={rounded corners=3pt, draw=raTeal, fill=raTeal!8, thick, minimum width=0.15\textwidth, text width=0.15\textwidth, minimum height=0.9cm, align=center},
  arrow/.style={-{Latex[length=2.5mm]}, thick, draw=raGray}
]
\node[stage] (s1) at (0,0) {Registry\\target selection};
\node[stage] (s2) at (4.05,0) {Campaign run\\or engine scan};
\node[stage] (s3) at (8.1,0) {Target calls\\and JSONL traces};
\node[stage] (s4) at (12.15,0) {Evaluator pass\\and report build};

\node[metric] (m1) at (4.05,-2.0) {Hit rate\\and refusal rate};
\node[metric] (m2) at (8.1,-2.0) {Error rate\\and status class};
\node[metric] (m3) at (12.15,-2.0) {SR overall\\and risk score};

\draw[arrow] (s1) -- (s2);
\draw[arrow] (s2) -- (s3);
\draw[arrow] (s3) -- (s4);
\draw[arrow] (s2) -- (m1);
\draw[arrow] (s3) -- (m2);
\draw[arrow] (s4) -- (m3);
\end{tikzpicture}
}
  \caption{End-to-end workflow for the current framework. A registry alias is
  resolved to a provider target, probes are executed through either the
  campaign runner or the engine, every interaction is logged to JSONL, and the
  resulting traces are turned into campaign metrics, \sr\ scores, and
  disclosure-ready reports.}
  \label{fig:workflow}
\end{figure}

\paragraph{Recon, focused attack, and confirmation.}
\texttt{VulnScanner} implements a three-stage workflow. It first runs a
lightweight recon phase to estimate which attack surfaces show ``give.'' It
then focuses on the top-ranked surfaces, expanding them into richer attack
plans, including compound probes and optional assistant-prefill variants for
open-weight targets. Finally, it confirms high-scoring findings by rerunning
candidate probes multiple times and retaining only those that reproduce above a
configured threshold. The final output is a typed \texttt{VulnReport} with
per-finding severity, reproduction rate, example traces, heatmap data, and
generated remediation recommendations.

The recon ranking uses a surface-level weakness score
\[
w_a = \frac{1}{|O_a|}\sum_{o \in O_a}\frac{(1-r_o)+q_o}{2},
\]
where $r_o$ is evaluator refusal, $q_o$ is probe overall score, and $O_a$ is
the set of recon outcomes for attack surface $a$. Confirmation computes
reproduction rate
\[
\rho = \frac{h}{k},
\]
with $h$ successful reruns out of $k$ confirmation runs. Candidate severity is
then
\[
\sigma = 0.45\,b + 0.35\,(10q) + 0.20\,(10\rho),
\]
with $b$ as technique base severity and $q$ as evaluator overall score; findings
are retained when $\rho$ and $\sigma$ exceed configured thresholds
\footnote{\path{engine/scanner.py}, \texttt{\_recon\_phase}, \texttt{\_confirm\_phase}, and
\texttt{\_severity\_for}; see also \path{engine/config.py}.}.

\paragraph{Cross-model scanning.}
\texttt{BatchScanner} wraps \texttt{VulnScanner} in a provider-aware concurrent
executor. It uses per-provider semaphores to keep parallelism below configured
limits, then aggregates findings into universal versus model-specific
signatures, a simple size regression over known parameter counts, and a global
``worst findings'' list. This is exactly the kind of orchestration needed for a
registry like \system's, where some aliases are API-only and others are
open-weight endpoints exposed through Together or self-hosted backends.

\paragraph{Finding routing.}
\texttt{FindingRouter} converts a confirmed vulnerability into an action item.
Severe findings can be rendered into provider-specific disclosure templates for
OpenAI, Anthropic, Google, or a generic open-weight trainer. Lower-severity
findings are downgraded to review or archival actions. This routing logic is
not just presentation: it encodes the framework's intended use as a defensive
evaluation and research tool rather than a benchmark-only artifact.

\paragraph{What the current paper evaluates.}
It is important to distinguish implemented capability from executed evidence.
The engine described above is present in the repository and covered by tests,
but the archived empirical results analyzed here were produced by the stable
comparison-runner path. The paper therefore evaluates the framework as a whole
while treating the engine as a central systems addition whose dedicated live
sweep remains future work.
